\documentclass{beamer}
\usepackage[utf8]{inputenc}
\usetheme{default}
\title{Binary JSON (BSON)}
\author{Arseny Kurnikov\\
Hasan Islam\\
Harri Hämäläinen}
\begin{document}


\begin{frame}
\titlepage
\end{frame}

\begin{frame}{Binary JSON}
    \begin{itemize}
        \item General purpose, binary encoded serialization format for JSON-ish data
        \item Supports all JSON types + binary
        \item Lightweight - minimum spatial overhead 
        \item Traversable - can skip some entries when parsing
        \item Efficient - quick encoding and decoding
        \item Easy to be implemented (partially) in most common environments
    \end{itemize}
\end{frame}

\begin{frame}{Lightweight}
    \begin{itemize}
        \item Minimize spatial overhead
        \item Adds overhead for faster traverse
        \item Human readability lost
        \item Saves up to 25\% space compared to JSON (though sometimes takes more than JSON)
    \end{itemize}
\end{frame}

\begin{frame}{Traversable}
    \begin{itemize}
        \item Length prefixes for elements
        \begin{itemize}
            \item \texttt{document ::= int32 e\_list "\textbackslash x00"}
            \item \texttt{string ::= int32 (byte*) "\textbackslash x00"}
        \end{itemize}
        \item Element types
        \begin{itemize}
            \item \texttt{element ::= "\textbackslash x01" e\_name double}
            \item \texttt{| "\textbackslash x02" e\_name string}
            \item \texttt{| "\textbackslash x03" e\_name document}
            \item \texttt{| "\textbackslash x04" e\_name document} (array)
            \item \texttt{| "\textbackslash x05" e\_name binary}
            \item \texttt{\dots}
            \item \texttt{e\_name ::= cstring}
        \end{itemize}
    \end{itemize}
\end{frame}

\begin{frame}{Efficient}
    \begin{itemize}
        \item C data types: \texttt{byte}, \texttt{int32}, \texttt{int64}
                            \texttt{double}
        \item Little-endian
        \item No need to parse text to encode/decode an integer or double
    \end{itemize}
\end{frame}

\begin{frame}[fragile]{Examples}
\begin{verbatim}
    {"hello": "world"}
    "\x16\x00\x00\x00\x02hello\x00
    \x06\x00\x00\x00world\x00\x00"
\end{verbatim}
    \begin{itemize}
        \item \texttt{\textbackslash x16\textbackslash x00\textbackslash x00\textbackslash x00}, length of the document is 22 bytes
        \item \texttt{\textbackslash x02}, string type (2)
        \item \texttt{\textbackslash x06\textbackslash x00\textbackslash x00\textbackslash x00}, the length of string is 6 bytes (including the terminating character)
    \end{itemize}

\begin{verbatim}
    {'a': 'foo', 'b': 1984}
    '\x17\x00\x00\x00\x02a\x00\x04\x00\x00\x00foo\x00
     \x10b\x00\xc0\x07\x00\x00\x00'
\end{verbatim}

    \begin{itemize}
        \item Little-endian, (0x07 $<<$ 8) + 0xc0 = 1984
    \end{itemize}
    
\end{frame}

\begin{frame}{Discussion}
    \begin{itemize}
        \item Based on JSON
        \item Readability: easily converted to JSON
        \item Suitable for both text and binary
        \item Bigger overhead in some cases
        \item Faster to parse
        \item Two string types (null terminated string with length prefix and null terminated cstring without length prefix)
    \end{itemize}
\end{frame}

\begin{frame}{Sources}
    \begin{itemize}
        \item BSON Specification, \url{http://bsonspec.org/}
    \end{itemize}
\end{frame}

\end{document}
